\honest{On Doing the Impossible}{Eliezer Yudkowsky}{2008}

High achievers give the advice ``Persevere.'' I used to think this meant working 14-hour days. I cannot do that; when work gets too hard I stop and read or watch something, and for years I thought I lacked the virtue. My younger self had the usual excuses: output counts, not input; laziness makes you back off failing methods; for creative work, peak output matters more than hours.

Then I looked back on my work in AI and saw that I ``had overestimated the difficulty of almost every single important problem.'' \nb{The post's evidence is that the problems came to seem less intimidating. It reports none solved, and later grants that ``you don't know how much work will be left when the confusion clears.''}

When I started, I pictured forty-year timescales and Manhattan Projects. \nb{Yudkowsky's 1996 essay ``Staring into the Singularity'' gives 2035, ``less than forty years!''} This is a common failure in AI futurism: the leap from ``I don't know how to solve this'' to imagining something big enough to feel as impressive as the problem. One man on the AI list says AI will cost a quadrillion dollars. Imagining something huge lets people think they can solve AI without understanding intelligence, and I made the same mistake. But having calculated that 55 million people die each year, 150,000 a day, I did not run away like a frightened rabbit. I asked what project could get there fastest; making the Singularity happen one hour earlier would repay a career. In 1998 I wrote a long treatise on ``seed'' AI, a term I coined, Brian Atkins, who had just sold Hypermart to Go2Net and later became the Singularity Institute's founding funder, asked whether a reasonable-sized team could actually do it. I told him it would take ``a Manhattan Project and thirty years,'' and for a while we considered a dot-com startup to raise the money. A year or two later, it seemed that a small organization could do preliminary work, such as new computer languages, and that open-source releases might be useful enough to grow, and the Singularity Institute started. That strategy was wrong, but the path looked shorter. \nb{The record agrees. ``Coding a Transhuman AI'' was announced in September 1998, and ``The Plan to Singularity'' (1999) aimed at a seed AI by 2010.}

I first ignored Friendly AI, since it seemed ``obviously impossible and useless'' to deceive a superintelligence about what is right; later I took it on as merely extremely difficult. I had also written off a precise understanding of intelligence as impossible, which removed it from my workload. That logic is deranged, since Nature does not care what you cannot do when It writes your requirements, but AI people still use it. Single problems came to seem smaller while the whole mountain grew, as problems moved from the impossible list to the to-do list.

Other AI researchers call ``Impossible!'' problems that seem to me ``eminently solvable.'' I name none. \nb{The link is to the Star Wars parody of four days before.}

``Impossible'' has two senses: a proof from stated axioms, and ``I can't see any way to do that.'' All my uses were the second. When you do not understand a domain, a query to your brain for a solution returns nothing. But there are mysterious questions, never mysterious answers. We know enough about stars to know one is hard to build, and enough about gears to prove that no gear train is a perpetual motion machine; those are bad problems for practicing the impossible. The confusing problems that feel most intimidating are where apparent difficulty is most likely to fall, since you do not know how much work will be left when the confusion clears. ``Confusion exists in the map, not in the territory.'' If you avoid blind alleys and have the ability, the ``apparent'' difficulty may fall. But a problem that seems impossible is not tried, a vicious cycle. Only because I was driven enough that ``forty years and a Manhattan Project'' meant starting sooner did I stay long enough to become less intimidated. Opposing biases rarely cancel, but here they did by luck: had I seen at the start that the task was a provably correct Friendly AI, not merely a seed AI, I would probably have burst into flames.

``By and large,'' AGI researchers are there because they think they have found ``the Key to AI'' that will work ``in just five years.'' \nb{The post gives no evidence for this. The post it links, ``Above-Average AI Scientists,'' rests on the author's impressions of AGI conferences.}

Richard Hamming asked colleagues what the important problems in their field were, and why they were not working on them. Important problems look big and scary; they promise no publications and no progress, perhaps for ten years. The most important problems are often impossible, which is why few philosophers work on reductionist accounts of consciousness. \nb{Hamming's talk says that what makes a problem important ``is that you have a reasonable attack.'' On his account, a problem with no visible way in is not yet an important one.}

This is not for everyone. Exceptional talent is only the ante; the chips are the years of your life, and ``you can lose.'' Most people should ``stick to the possible.'' Never give up? ``Don't be ridiculous.'' Knowing when to lose hope is a skill. Perseverance keeps things difficult, and there are easier ways to get glamour and respect. But if something is worse to you than wasting your life, you may have cause to try.

Because staying at work was a constant struggle, it was what I noticed about myself; I did not see that perseverance applies at other timescales until I saw people declare ``impossible'' at once anything they did not want to try, or shy from work that might take decades instead of five years. On the scale of seconds, perseverance is not giving up at the first sign of difficulty; on the scale of years, it is staying with an insanely hard problem when you could get better rewards elsewhere. So it works at three timescales: not running away takes seconds, working takes hours, sticking at it takes years. The first I had to learn; the second is still a struggle; the third comes naturally.
